% !TeX root = ../colloquium-2.tex

% Свойство линейности неопределённого интеграла. 
% Формулировка теоремы о формуле интегрирования по частям.
% Формулировка теоремы о замене переменной в неопределённом интервале.

\phantomsection\label{sec_Definition21}
\begin{definition}[
    Свойство линейности неопределённого интеграла. 
    Формулировка теоремы о формуле интегрирования по частям.
    Формулировка теоремы о замене переменной в неопределённом интервале.
]
    \subsection*{Свойство линейности неопределённого интеграла. }
    \begin{note}[purple]
        Если функции $F$ и $G$ дифференцируемы на интервале $(a, \; b)$ и $F' = f$,
        а $G' = g$, то в силу линейности производной при $\alpha, \beta \in \mathbb{R}$
        $$ (\alpha F + \beta G)' = \alpha f + \beta g $$
    \end{note}

    \subsection*{Формулировка теоремы о формуле интегрирования по частям.}
    \begin{note}[purple]
        Пусть функции $f$ и $g$ дифференцируемы на интервале $(a, \; b)$. Пусть также
        на этом интервале существует первообразная функции $f'g$. Тогда на интервале $(a, \; b)$
        существует первообразная функции $g'f$ и справедлива \textbf{формула интегрирования 
        по частям}. 
        $$ \int f'(x)g(x)dx = f(x)g(x) - \int f(x)g'(x)dx + C $$
    \end{note}

    \subsection*{Формулировка теоремы о замене переменной в неопределённом интервале.}
    \begin{note}[purple]
        Пусть функция $\varphi$ определена и дифференцируема на интервале $(\alpha, \; \beta)$
        и при любом $t \in (\alpha, \; \beta)$ значение $\varphi(t)$ принадлежит интервалу
        $(a, \; b)$, на котором определена функция $f$. Пусть также существует интеграл 
        $\int f(x)dx$. Тогда существует интеграл $\int f(\varphi(t)) \varphi'(t) dt$ и
        верна следующая формула замены переменной:
        $$ \int f(x)dx = \int f(\varphi(t)) \varphi'(t) dt $$
    \end{note}
\end{definition}